\clearpage\thispagestyle{empty}
\pdfbookmark[0]{Quellen zu Gauß an Bode (7. September 1818)}{bodesep1818-sources}
\begin{center}\textbf{Quellen zu Gauß an Bode (7. September 1818)}\end{center}
\noindent Diese KI-Übertragung ins heutige Deutsch folgt Carl Friedrich Gauß’ Brief an Johann Elert Bode in den \href{https://gdz.sub.uni-goettingen.de/id/PPN236007467}{\textit{Werken}, Band~VI} (Göttingen, 1874), S.~414--416. Der Erstabdruck erschien unter dem Titel \textit{Ueber den Repsoldschen Meridiankreis, Beobachtung des Uranus, Saturns und der Pallas etc.} als Beitrag~24 in Bodes \href{https://api.digitale-sammlungen.de/iiif/image/v2/bsb10538350_00003/full/1600,/0/default.jpg}{\textit{Astronomisches Jahrbuch für das Jahr 1821}} (Berlin, 1818), S.~\href{https://api.digitale-sammlungen.de/iiif/image/v2/bsb10538350_00218/full/1600,/0/default.jpg}{212}--\href{https://api.digitale-sammlungen.de/iiif/image/v2/bsb10538350_00222/full/1600,/0/default.jpg}{216}. Die Links zum Jahrbuch führen zum Exemplar der Bayerischen Staatsbibliothek.
\par\smallskip\noindent Im Erstabdruck stehen auf S.~\href{https://api.digitale-sammlungen.de/iiif/image/v2/bsb10538350_00219/full/1600,/0/default.jpg}{213} zwei Anmerkungen. Die erste gehört zur Erwähnung der Pallas-Ephemeride: „Folgt weiter unten.“ Die von Dirksen berechnete Ephemeride erschien gesondert als Beitrag~30, S.~\href{https://api.digitale-sammlungen.de/iiif/image/v2/bsb10538350_00230/full/1600,/0/default.jpg}{224}--\href{https://api.digitale-sammlungen.de/iiif/image/v2/bsb10538350_00231/full/1600,/0/default.jpg}{225}. Sie reicht vom 7. August 1819 bis zum 6. Juni 1820 und gilt für Mitternacht in Göttingen. Die zweite Anmerkung gehört zu Gauß’ Bemerkung, dass Bode den Kreis zweifellos schon gesehen habe: „Der geniale Instrumentenbauer zeigte ihn mir, als ich im Jahr 1810 in Hamburg war. B.“
\par\smallskip\noindent Im Erstabdruck steht auf S.~\href{https://api.digitale-sammlungen.de/iiif/image/v2/bsb10538350_00220/full/1600,/0/default.jpg}{214} bei der ersten Gruppe der Polarsternbeobachtungen \textit{Jul.}, in der Werkausgabe steht \textit{Juni}. Auf S.~\href{https://api.digitale-sammlungen.de/iiif/image/v2/bsb10538350_00221/full/1600,/0/default.jpg}{215} hat der Mittelwert der Korrektur einen Doppelstrich, in der Werkausgabe ein hochgestelltes \textit{s}. Die Deklination des Uranus am 3. Juni endet im Erstabdruck auf $41'',5$, im Brief in der Werkausgabe auf $41''6$. Die Verbform „drücken“ nach der Einzahl „Spalte“ und die Schreibung „Boobachtungen“ folgen der Werkausgabe. KI-Übertragung und KI-Satz.
\clearpage
